% Part: computability
% Chapter: computability-theory
% Section: non-comp-set

\documentclass[../../../include/open-logic-section]{subfiles}

\begin{document}

\olfileid{cmp}{thy}{ncp}
\olsection{محاسبه نۀ کېدونکي سټونه شته}


پورته مو وليدل چې هر محاسبه کېدونکے سټ په محاسبوي ډول د شمېر وړ
دے۔ ايا د دې معکوس هم رښتيا دے؟ لاندې قضيه ښيي چې په عمومي ډول داسې
نۀ ده۔

\begin{thm}\ollabel{thm:K-0}
فرض کړئ $K_0$ سټ $\Setabs{\tuple{e, x}}{\cfind{e}(x) \fdefined}$
دے۔ بيا $K_0$ په محاسبوي ډول د شمېر وړ دے، خو محاسبه کېدونکے نۀ دے۔
\end{thm}

\begin{proof}
د دې د ليدلو دپاره چې $K_0$ په محاسبوي ډول د شمېر وړ دے، پام وکړئ
چې دا د هغې تابع~$f$ د تعريف ساحه ده چې داسې تعريفېږي:
\[
f(z) = \umin{y}{(\len{z} = 2 \land T((z)_0, (z)_1, y))}.
\]
ځکه که $\cfind{e}(x)$ تعريف شوې وي، $f(\tuple{e, x})$ د درېدلې
محاسبې يوه لړۍ مومي؛ که $\cfind{e}(x)$ تعريف شوې نۀ وي، نو
$f(\tuple{e, x})$ هم تعريف شوې نۀ ده؛ او که~$z$ ان يوه جوړه هم نۀ
کوډوي، نو $f(z)$ بيا هم تعريف شوې نۀ ده۔

دا حقيقت چې $K_0$ محاسبه کېدونکے نۀ دے، هماغه د درېدنې د مسئلې
ناپرېکړتيا ده، \olref[hlt]{thm:halting-problem}۔
\end{proof}

سټ~$K_0$ د هغو جوړو~$\tuple{e,x}$ سټ دے چې
$\cfind{e}(x) \fdefined$ وي، يعنې $\tuple{e,x} \in K_0$ هله او
يوازې هله چې~$\cfind{e}$ پر ننوت~$x$ تعريف شوې وي (ودرېږي)، نو دې
ته ``د درېدنې سټ'' هم ويل کېږي۔ سټ
$K = \Setabs{e}{\cfind{e}(e) \fdefined}$ د ``پر خپل ځان د درېدنې
سټ'' دے۔ دا ډېر ځله د يوې معياري ناپرېکړه کېدونکې بېلګې په توګه
کارول کېږي۔

\begin{thm}\ollabel{thm:K}
پر خپل ځان د درېدنې سټ
$K = \Setabs{e}{\cfind{e}(e) \fdefined}$، !!{c.e.} دے خو د پرېکړې
وړ نۀ دے۔
\end{thm}

\begin{proof}
  فرض کړئ $K$ د پرېکړې وړ دے، يعنې ځانګړونکې تابع~$\Char{K}$ ئې
  محاسبه کېدونکې ده۔ وټاکئ:
  \[d(e) = \begin{cases}
  1 & \text{که $\Char{K}(e) = 0$ وي}\\
  \fundefined & \text{که داسې نۀ وي۔}
  \end{cases}
  \]
  فرض کړئ $k$ د~$d$ شاخص دے، يعنې $d \simeq \cfind{k}$۔ بيا
  $d(k) \simeq \cfind{k}(k)$ ده۔ دا له هغه حقيقت سره تناقض لري چې
  $d(k) \fdefined$ هله او يوازې هله چې
  $\cfind{k}(k) \fundefined$ وي، او دا د~$d$ له تعريفه راځي۔

  $K$ د تابع $f(x) = \umin{y}{T(x,x,y)}$ د تعريف ساحه ده او ځکه
  !!{c.e.} دے۔
\end{proof}

\end{document}
